\section{Introduction: Panpsychism's Appeal}

Panpsychism derives much of its contemporary appeal from an epistemic asymmetry. Consciousness, unlike many entities posited by physical science, is known from the first-person point of view. When I am in pain, I do not posit the pain as I posit an electron or a gene; the pain is given. Physical science, by contrast, is often taken to describe matter in structural, relational, and dispositional terms. It tells us what matter does, but not, at least on one influential line of thought, what matter is intrinsically. If consciousness is the one reality known from the inside, then perhaps consciousness, or something consciousness-like, belongs in our account of fundamental reality. This thought gives panpsychism its distinctive force. It promises to avoid the apparent emergence of consciousness from an entirely non-conscious base by locating mentality, proto-mentality, or phenomenal character in the nature of matter itself. In its more sophisticated forms, the view is not the simple claim that electrons have recognizably human-like minds. Strawson's realistic monism, for example, argues that a serious physicalism must take experience as part of the physical world and cannot treat experiential reality as an inexplicable late arrival \citep{strawson2006}. Goff's Russellian monist defense places panpsychism in a broader argument about the categorical nature of matter and the limits of structural physics \citep{goff2017cfr}. The field also divides between micropsychist views, which begin from many small subjects and try to combine them into human minds, and cosmopsychist views, which begin from a single cosmic subject and try to derive local minds from it; both forms figure in what follows.

This paper isolates one prominent motivation for panpsychism: the claim that consciousness has a special epistemic status because it is known directly from within. Panpsychists may also appeal to the alleged impossibility of radical emergence, to continuity in nature, to theoretical simplicity, to dissatisfaction with standard physicalism, to the explanatory gap, or to a Russellian monist thesis according to which physics leaves the intrinsic nature of matter unspecified. Those arguments raise large issues that cannot be settled here.\footnote{This bracketing leaves important arguments in place. Alter develops a recent route from the knowledge argument to Russellian monism, Pereboom treats Russellian monism as part of the prospects for physicalism, and Stoljar's ignorance-based treatment warns against moving too quickly from epistemic limitation to metaphysical conclusion \citep{alter2023matter,pereboom2011physicalism,stoljar2006ignorance}.} The narrower question is how much work the first-person beginning can do once the view introduces micro-experience, proto-experience, cosmic subjectivity, or other hidden phenomenal posits.

Panpsychism remains a live metaphysical option, but one of its most powerful motivations is weaker than it looks. If the hidden phenomenal base is not itself disclosed in experience, panpsychism has not dissolved the emergence problem in advance. The gap reappears between that theoretical base and the consciousness we actually undergo. Goff's recent treatment of subjects as strongly emergent makes this reappearance visible from within his own framework \citep{goff2024}.

\section{The Datum}

First-person experience gives panpsychism a strong beginning. It establishes that consciousness is real, and it discloses something about what conscious life is like. This is not a small concession. A theory of consciousness that denies or deflates the first-person datum carries a substantial burden.

Three questions must be kept apart: whether conscious experience exists, what features it has, and what its metaphysical status is. First-person experience answers the first robustly. It bears on the second more cautiously, since our grasp of phenomenal character may be partial, vague, and fallible. It does not by itself answer the third. Acquaintance with experience does not settle whether consciousness is fundamental, intrinsic, derivative, organizational, microphysical, or cosmic.

The datum supplied by experience is real but finite. It establishes consciousness at the grain of acquaintance: experience as given, with some phenomenal character however attenuated. A mood may be vague, peripheral vision indeterminate, and anticipation foggy; these are still ways experience is given. Vagueness and fogginess are phenomenal characters, not routes to subjectivity without character. The same point applies to harder cases, including minimal experience, meditative states, moods, objectless anxiety, dreamless-sleep reports, and claims about pure consciousness. Those cases can be granted. What they grant is still experience as given, not a metaphysical verdict about the distribution, fundamentality, or cosmic structure of consciousness.

The live question is therefore one of reach. First-person experience is indispensable evidence for consciousness. It does not automatically support claims about micro-experience, proto-experience, proto-phenomenal categorical bases, cosmic subjectivity, or phenomenal bonding relations.

\section{The Constraint}

\subsection{The constraint stated}

Panpsychism often draws force from a claim about the special epistemic status of consciousness. Consciousness is not known in the way electrons, genes, or black holes are known. I do not posit my current pain as a theoretical entity. I undergo it. Conscious experience is given from the first-person point of view.

Direct knowledge of consciousness gives warrant for what is actually given in experience. It does not give every entity described as phenomenal, proto-phenomenal, or experiential the same standing. First-person acquaintance gives immediate warrant for the occurrence and character of the experience one is acquainted with. It does not, by itself, warrant claims about the distribution, fundamentality, categorical role, or cosmic or microphysical basis of consciousness. That is the introspective warrant constraint:

\begin{quote}
A theory that motivates itself by appeal to first-person experience may claim first-person warrant only for what first-person experience actually discloses.
\end{quote}

The constraint is modest. It allows theoretical posits. Physics, biology, and psychology all legitimately posit entities and structures not given in experience. Those posits are warranted by explanatory, predictive, or unificatory work. They do not inherit the authority of first-person acquaintance. If panpsychism posits micro-experience, proto-experience, cosmic subjectivity, hidden phenomenal fields, or relations of phenomenal inheritance, those posits may be defended as theoretical entities. They may not be treated as if they were directly warranted by the first-person datum of consciousness.

Put as an argument:

\begin{enumerate}
\item First-person acquaintance warrants only what experience actually discloses.
\item What experience discloses is consciousness at the grain of acquaintance: experience as given, with some phenomenal character however attenuated.
\item It does not disclose micro-experience, proto-experience, proto-phenomenal categorical bases, cosmic subjectivity, or relations of phenomenal bonding or inheritance.
\item So those posits do not inherit first-person warrant.
\item If they are defended at all, they must be defended as theory, on the same abductive terms as any other metaphysical hypothesis.
\end{enumerate}

The first step concerns warrant, not existence. It does not say that only what experience discloses is real, and it leaves every panpsychist posit metaphysically open. It says that such posits are not underwritten by first-person acquaintance, and so must find their support elsewhere. A panpsychist may reply, correctly, that no careful defender claims first-person experience alone proves panpsychism. The stronger argument is abductive: consciousness is known directly; physics gives only structure; radical emergence looks suspect; panpsychism is offered as the best metaphysical hypothesis. The remaining problem concerns the weight assigned to the first-person datum. Even framed abductively, that datum can be made to carry more weight than the abduction earns.

\subsection{Why hidden posits do not inherit first-person warrant}

The constraint becomes most important when panpsychism moves from consciousness as given in acquaintance to hidden phenomenal posits. Contemporary panpsychist and panprotopsychist views may posit micro-experience, proto-experience, proto-phenomenal properties, cosmic subjectivity, hidden phenomenal fields, phenomenal bonding relations, or forms of phenomenal inheritance. These posits are not directly disclosed in first-person experience. The immediacy of pain warrants the pain, not a cosmic subject behind it; the directness of visual experience warrants that visual experience occurs, not that the microphysical entities involved in vision instantiate hidden phenomenal natures.

Revelation is the strongest principle by which acquaintance might be thought to reach farther. In its strong form, defended by Goff, phenomenal concepts ``reveal the complete nature of the conscious states they refer to'': attending to a pain or a color experience, one grasps what that experience is \citep[106--132]{goff2017cfr}. The constraint runs even on that strong version, for a precise reason. Grant that introspection reveals the complete nature of the phenomenal quality one undergoes. The quality is one thing; the metaphysical role of the property that bears it is another. Acquaintance may reveal what a phenomenal quality is like. It does not reveal that the property so encountered plays the categorical-grounding role. That role assignment is theoretical. Roelofs's moderate revelation thesis makes the same limit visible from the other direction: constitutive panpsychism can be protected only if introspection need not disclose every constitutive fact transparently \citep{roelofs2020revelation}. That concession is not a defeat for panpsychism. It marks the point at which hidden constitution must be defended by theory rather than by the direct authority of acquaintance.

The warrant issue intersects the combination literature.\footnote{The nearby literature is extensive. Bruntrup maps the contemporary combination-problem terrain; Chalmers distinguishes panpsychism from panprotopsychism and treats combination as a central obstacle; Coleman presses the subject-summing problem; M{\o}rch develops a phenomenal-powers view; Roelofs defends composite subjectivity; Nagel's classic discussion asks whether panpsychism follows from claims about mental properties, reduction, and composition \citep{bruntrup2017,chalmers2015,chalmers2017combination,coleman2014,morch2018,morchForthcoming,roelofs2019,nagel1979panpsychism}.} Panpsychism's largest internal puzzle is how many small minds, micro-experiences, or proto-phenomenal bases could yield the macro-consciousness we know from our own case. The formulation already contains posits not delivered by the first-person datum: micro-subjects, micro-experiences, bonding relations, structural correspondences between micro- and macro-experience, or powers whose phenomenal nature grounds physical dispositions. If those posits solve the combination problem, they may be justified by that success. Such justification would be theoretical. It would not be the direct authority of first-person acquaintance carried downward or outward.

Panprotopsychism makes the point especially clear. The panprotopsychist need not say that fundamental entities are conscious. The claim is that they instantiate proto-phenomenal or proto-experiential properties that can help explain consciousness. But proto-experience is not something we encounter from within. Introspection is, by what it does, acquaintance with experience; its reach is whatever reach experience has. Calling something proto-experiential in the same breath as experience does not extend acquaintance to it. If proto-experience is sufficiently like experience, it inherits familiar problems about structure, unity, subjecthood, and combination. If it is sufficiently unlike experience, it loses its claim to first-person warrant.

Phenomenal bonding faces the same limit. Goff's discussion of phenomenal bonding is explicitly designed to address the combination problem, the problem of how many micro-level subjects or phenomenal properties could yield one macro-level subject \citep{goff2017bonding}. That machinery may be part of a serious metaphysical theory. It is not given in first-person experience. I am acquainted with pain, color, thought, and temporal flow. I am not acquainted, simply by being conscious, with relations by which micro-subjects combine into macro-subjects.

\subsection{How the constraint bears on specific views}

The constraint bears differently on different panpsychist views. Some treat first-person experience as directly supporting hidden phenomenal ontology; the criticism applies to them most directly. Others argue abductively in their official statements while still drawing dialectical force from first-person privilege, and these are the central case here. Others make no first-person warrant claim at all and rest the view entirely on theoretical virtues; the criticism leaves them largely untouched. Saad's argument from psychophysical harmony is the clean example of this third kind \citep{saad2022}.

The worry reaches careful formulations as well as careless ones. Strawson's realistic monism is the austere baseline. He argues that experience is real, that experience is physical if physicalism is true, and that the emergence of experience from the wholly non-experiential would be unintelligible \citep{strawson2006}. His text leaves the move on the surface. Experience is ``the fundamental given natural fact,'' and for that reason ``the obligatory starting point for any remotely realistic'' theory of what there is \citep[4]{strawson2006}. The inference runs from the certainty of the datum to a constraint on fundamental theory: nothing is more certain than that experience occurs, so any realistic account of what exists must build experience in at the ground. Strawson may be right that experience is part of reality and cannot be dismissed. First-person acquaintance with experience still does not warrant the stronger claim that experientiality exists at the fundamental level.

Goff's more elaborate Russellian and cosmopsychist machinery makes the problem sharper. His work develops panpsychism within a structure/dynamics framework and addresses the combination problem through phenomenal bonding and cosmopsychism \citep{goff2017cfr,goff2017bonding,goff2019cosmopsychism}. His Revelation thesis carries much of the weight in his case against physicalism \citep[106--132]{goff2017cfr}, and the same quality/role distinction applies: whatever acquaintance reveals about the character of an experience, it does not reveal the metaphysical role of that experience. His positive case then turns on a simplicity argument. The intrinsic nature of the matter of brains is taken to be consciousness-involving, and the most parsimonious hypothesis is that matter elsewhere is continuous with it. As an abduction this is legitimate, and it claims no acquaintance with electrons; the pressure falls instead on its starting point. That the intrinsic categorical nature of brain matter is consciousness-involving is a metaphysical reading of acquaintance, not a deliverance of it. Acquaintance gives the experience, not its identification with the categorical nature of the matter.

Cosmopsychism reverses the micropsychist direction. Instead of beginning with many micro-subjects and combining them into human subjects, it begins with a cosmic subject and derives local subjects from it. This may avoid some versions of the combination problem, but it intensifies the introspective warrant problem. Cosmic subjectivity is hidden from third-person science and unlike anything first-person experience discloses. First-person experience reveals local, finite, perspectival consciousness with determinate phenomenal character. It is also bounded and affectively grained: one is acquainted with one's own experiences and not with anyone else's, and what is given is given as mattering in some way, pleasant or aversive or neutral. Perspectival closure and affective valence both belong to the datum, and neither is delivered by a unified cosmic experiencer from which local minds derive.

That claim should not be read as assuming that cosmopsychism faces a simple mirror image of micropsychism's combination problem.\footnote{Shani resists that framing, arguing that cosmopsychic individuation need not be understood as the reverse of micropsychist subject-summing. Albahari and M{\o}rch likewise explore versions of cosmic consciousness that complicate simple charges about decombination and unity \citep{shani2015,shani2022,albahari2019,morch2024allone}.} Even if these views offer a coherent theory of cosmic individuation, the cosmic subject is not disclosed in first-person acquaintance. Its warrant is holistic, metaphysical, or explanatory, not directly first-personal.

Goff's recent hybrid cosmopsychism changes the target again. He now treats subjects as strongly emergent while phenomenal properties are weakly emergent \citep{goff2024}. The division is meant to let panpsychism help with consciousness without inheriting the worst of the subject-combination problem. Goff's view need not collapse for that concession to matter. The anti-emergentist motivation has narrowed. It no longer covers subjecthood itself, and so cannot by itself secure the full first-person datum. Once panpsychism explains subjecthood through emergence, grounding, and top-down structure, it is doing theoretical metaphysics, and cannot retain the immediacy of the original first-person datum for the theoretical phenomenal base it introduces.

M{\o}rch's phenomenal powers view is a partial case. On that view, phenomenal properties are not idle qualia but powers that can ground physical dispositions \citep{morch2018,morchForthcoming}. If successful, this would give panpsychism a substantive explanatory role. It also makes the constraint sharper. First-person experience may acquaint us with phenomenal character, and perhaps with agency, motivation, or felt efficacy in ordinary cases. It does not follow that we are acquainted with fundamental phenomenal powers that ground the dispositional structure described by physics. M{\o}rch's argument is explicitly theoretical, and to that extent not the paper's target. But it is motivated by a knowledge claim: the case for phenomenal grounds runs through the premise that ``the only categorical properties we know are phenomenal properties'' \citep{morch2018}. The explicit argument is abductive, while the consideration that recommends it is the privilege the constraint examines.

A constitutive panpsychist may press the point in the other direction, and Roelofs develops the most careful version: if the consciousness we undergo is literally constituted by micro-experiences, then introspecting it already amounts to indirect acquaintance with what makes it up. Constitutive panpsychism, Roelofs argues, need not predict introspective access to micro-properties: they constitute what is introspected without themselves appearing as objects of introspection. He grants that we are ``unable to tell introspectively just how composite'' our experience ``might really be'' \citep[89]{roelofs2020revelation}. The reply is that constitution and acquaintance are different relations. Cells constitute me, but I am not acquainted with my cells by being acquainted with myself. Even granting the constitutive picture, introspection presents the integrated whole rather than the items that compose it. Acquaintance presents at some grain, and that is the grain at which warrant is given; a coarser-grain acquaintance does not by itself deliver finer-grain acquaintance with whatever constitutes it. Roelofs warns against ``inferring absence of structure from the failure of structure to be manifest'' \citep[87]{roelofs2020revelation}, and the warning is well taken, since the constraint draws no such inference. Its conclusion is about warrant, not existence: introspection does not warrant the micro-level, and that leaves its existence entirely open. Roelofs's own account agrees, treating the micro-level as introspectively undisclosed and defensible only by its theoretical role in explaining the macro-consciousness that is disclosed.\footnote{Lin argues that revelation cannot favor panpsychism over physicalism, while Roelofs argues that constitutive panpsychists can accept a medium-strength revelation thesis without transparent introspective access to every constitutive fact. Ramm gives the present worry a more direct target by using phenomenological reflection to motivate an idealist-panpsychist picture \citep{lin2025,roelofs2020revelation,ramm2021}.}

The constraint targets a pattern of warrant, not individual writers. When a panpsychist treats the inference as fully abductive, the constraint records that fact rather than objecting to it. The comparative case for panpsychism often depends on the claim that it takes consciousness more seriously than rival views can. Once the micro-level, cosmic level, or bonding relation is not introspectively disclosed, first-person evidence no longer warrants it directly. The view may still be defended. Its support is then theoretical.

\section{The Abstraction Problem}

The first-person datum is sometimes described as bare subjectivity, pure awareness, or mere what-it-is-likeness. That description is tempting because it seems to isolate the minimal feature common to all conscious states. It also risks abstracting away the features that make consciousness intelligible as a datum.

When we attend to consciousness, we do not find subjectivity as an empty item standing apart from all character. In familiar waking cases, we find sensation, mood, thought, temporal flow, embodiment, and varying degrees of attention. Less familiar cases should be treated with care. A meditative state, an objectless mood, or a claimed pure-consciousness state may lack many features of standard waking experience. But insofar as it is available to first-person acquaintance, it is available as something: calm, blank, diffuse, silent, contentless, expansive, or otherwise phenomenally determinate. The relevant point is not that every conscious state represents an external object. It is that the experiences we can appeal to as evidence are never given as bare metaphysical subjectivity.

Appeals to fundamental subjectivity now face a dilemma. Does the posited fundamental consciousness have phenomenal character, or not? If it does, then it already has structure of some determinate sort. If it does not, then it is not something first-person experience discloses at all, but an abstraction from experience rather than an item found in it.

The dilemma is especially important for views that retreat from rich consciousness to minimal subjectivity. The move is understandable. Few panpsychists want to say that electrons have visual fields, memories, anxieties, or personal identities. So the relevant posit is often thinned out. The fundamental level is said to contain minimal phenomenal character, bare subjectivity, or proto-experiential nature. But thinning the posit changes its epistemic status. The thinner it becomes, the less it resembles the consciousness that first-person experience actually gives. Thinning also does not escape the first horn of the dilemma. Whatever phenomenal character is retained, however minimal, is character of some determinate sort; structure does not vanish by being made small. A posit with minimal phenomenal character still has phenomenal character, and first-person acquaintance with our experience does not by itself warrant that character's distribution, fundamentality, or microphysical placement.

A natural reply is that panpsychism is doing what scientific theorizing often does: abstracting a property from familiar cases and assigning it to entities in which it does not appear in familiar form. We abstract mass from heavy bodies and assign it to particles we cannot feel. Why not abstract phenomenality from experience and assign it to the fundamental level?

The analogy is useful, but limited. Mass is embedded in a network of measurement and intervention that supplies warrant for applying it beyond the cases in which we first encounter it. We know what it would be for an unfamiliar entity to have mass, because mass has behavioral signatures that can be tracked without acquaintance. Its identity is fixed by its role in theory and practice, not by the felt heaviness of familiar objects. Phenomenality is different. We do not encounter it apart from there being some way the experience is given. If phenomenality is stripped of temporality, perspective, valence, and every determinate mode of givenness, it is no longer clear what property has been isolated. The panpsychist may answer that the residue is simply what-it-is-likeness. But naming the residue does not escape the dilemma: either it is some determinate way an experience is given, and so has phenomenal character, or it is no way of being given at all. Mass avoids this fork because its identity never rested on a felt mode of givenness. For phenomenality the determinate modes of givenness do the individuating work, and the abstraction is precisely their removal. The panpsychist may still posit a thin phenomenal property, but the posit is now theoretical. It no longer carries the first-person authority of the experience from which the abstraction began.

A related reply says that minimal phenomenal character can be primitive at the fundamental level, in the sense that it is not further analyzable into more basic phenomenal constituents, while still being determinate. The first horn of the dilemma is then conceded, but its consequence is denied: determinate character need not call for reductive explanation if no further phenomenal analysis is available. The primitive can be granted without granting the warrant transfer. Acquaintance with determinate character in our own case does not by itself license claims about where such character is distributed, whether it is fundamental, or whether it belongs at the microphysical level. Granting the primitive's coherence also does not yet pay the debt that matters most to the panpsychist: the route from a primitive micro-character to the determinate macro-consciousness we have is still owed. The dilemma's purpose was not to require reductive analysis of phenomenality. It was to show that bare characterless subjectivity cannot be the datum, and that any determinate alternative cannot simply inherit the first-person warrant of the experience from which it was abstracted.

This point overlaps with Frankish's worry about the depsychologization of consciousness in panpsychist contexts \citep{frankish2021}. The more consciousness is stripped of the psychological roles by which we ordinarily recognize it, the less clear it becomes that the retained posit is still the consciousness from which the argument began. I do not rely on an illusionist conclusion. The more modest point is enough: a depsychologized or contentless posit may be coherent, but it is not simply the first-person datum under another name.

Even at the most attenuated end of waking consciousness, what is given is not nothing. There is some sense of presence, of for-me-ness, of the moment's being a moment for the one undergoing it. That sense is itself a phenomenal character, however thin, and not a structureless occupant of consciousness. This is a phenomenological claim about what acquaintance presents, not a claim about what a later report or memory can reconstruct.

\subsection{Pure Awareness}

The strongest challenge here comes from meditators who report pure or contentless awareness: moments of consciousness without object, self, or thought, exactly the bare subjectivity the preceding argument denies is given in experience. The challenge can be granted without granting its conclusion. If pure awareness has some phenomenal character, even an attenuated one such as silence, openness, presence, or bare temporality, it is not characterless subjectivity. If it has no phenomenal character at all, it is not disclosed as experience in the way pain or temporal flow is disclosed. Either way, the appeal to pure awareness does not license the transfer from first-person acquaintance to a hidden ontology of bare fundamental consciousness.

The literature gives the challenge real weight. Metzinger develops minimal phenomenal experience as a candidate for the simplest form of consciousness, drawing on meditative reports of pure awareness and tonic alertness \citep{metzinger2020mpe}; Gamma and Metzinger, surveying meditators on what is ``often described as a contentless form of experience,'' find that such reports still resolve into a multi-factor phenomenal profile, with dimensions such as felt time, silence, peace, and wakeful presence \citep{gammaMetzinger2021}; and philosophical treatments of witness-consciousness and pure awareness suggest that awareness may be present without object-directed content, self-identification, or discursive thought \citep{albahari2006,albahari2009,thompson2014}. Meditation-induced cessations make the evidential problem vivid: what reaches the discussion is always a report of an episode in which ordinary forms of consciousness were said to be absent \citep{laukkonen2023}.

These cases should be granted for the sake of argument. I need not deny that there are minimal, nondual, selfless, or apparently contentless states, nor decide whether nirodha samapatti is best described as consciousness without content, absence of consciousness, or a remembered transition around an absence. The evidential issue arises when such states are used in public philosophical argument. They reach the theorist through report, memory, interpretation, and later classification, and a later report does not settle what structure the episode had. A sympathetic reader may object that the report is only a later trace, while the evidence itself is the private episode. That distinction matters. Private acquaintance, if it occurs, may have first-person force for the subject before any report is made. Once the episode is offered as a public reason for panpsychism, however, it enters the argument through memory, description, interpretation, and classification. Even apart from that public step, the appeal still divides the same way: if the episode was given as silence, openness, presence, bare temporality, or some other way of seeming, it had phenomenal character; if nothing was given, there is no first-person disclosure of experience for panpsychism to extend.

A report used as evidence carries built-in commitments. A reportable episode must be retained or reconstructed in memory and attributed to the subject who reports it; Lin presses this for selfless states, where it is hard to see how an autobiographical memory could establish what a wholly selfless episode was like \citep{lin2022selfless}. It must also be located, at least minimally, in a before-and-after order; Windt argues that even content-poor sleep experience involves a minimal temporality, a duration or a sense of having been \citep{windt2015}. And experiential subjectivity may not detach cleanly from a minimal for-me-ness or bodily perspective at all, as Zahavi argues and as Henry and Thompson press against Albahari \citep{zahavi2005,henryThompson2011}. Each point cuts the same way. If the structure is present, the report attests to more than bare subjectivity. If it is absent, the retrospective report cannot establish the episode as a first-person datum lacking it.

Contentless-awareness reports still have evidential value. At most, they support attenuated phenomenal character, reduced object-directedness, weakened self-representation, or unusual temporal and bodily structure. The more a report says was absent during the episode, whether memory, selfhood, or temporal grain, the more cautiously it must be used, since the ordinary conditions for accurate retrospective report are exactly what it says were suspended. Schwitzgebel's worries about introspective unreliability compound the restriction rather than replacing it \citep{schwitzgebel2008,schwitzgebel2011}.

\section{The Russellian Vacancy}

The Russellian monist can press a further argument. Physics, it is said, gives us only structure and dynamics. It describes relations, dispositions, equations, and causal roles. It leaves open the intrinsic or categorical nature of the entities that bear those roles. Consciousness, by contrast, is known from within. Perhaps consciousness reveals what matter is intrinsically, or at least gives us a model for what intrinsic nature could be \citep{goff2017cfr,alter2023matter,pereboom2011physicalism}.

The argument begins from a real limitation in structural description rather than a simple appeal to mystery. If a theory tells us only how entities are related or how they behave, one may ask what, if anything, has the structure or grounds the dispositions. Alter's recent defense of Russellian monism makes the vacancy more than a rhetorical device: the claim is that standard physical description omits intrinsic properties that may both underlie physical structure and bear on consciousness \citep{alter2023matter}. Pereboom's discussion matters for a different reason. It shows that Russellian monism need not be introduced as a straightforward rejection of physicalism, but can be explored as part of a physicalist response to arguments from consciousness \citep{pereboom2011physicalism}. These uses differ, but they share the thought that physical description may leave an intrinsic-nature question open.

Two claims should be kept distinct. One is that physics leaves categorical nature unspecified. The other is that consciousness is our only acquaintance with categorical nature. The first claim identifies a vacancy. The second tries to fill it. The pressure falls on the second claim, because first-person experience gives consciousness as phenomenal character, not as categorical nature already interpreted. We know pain as pain, color-experience as color-experience, mood as mood. Whether the property so encountered is intrinsic categorical nature, derivative organization, emergent process, or something else is a further metaphysical reading of the datum.

The pull toward phenomenal filling is understandable. If physics leaves categorical nature unspecified, and if consciousness is the one concrete reality whose nature seems available from the inside, then consciousness appears to be the only available model of intrinsic nature. The move is not arbitrary; it tries to use the one apparent case of acquaintance to illuminate what structural physics leaves dark. But the fact that the model is the only one available tells us about the reach of our concepts rather than the nature of matter. A vacancy we cannot characterize is better left uncharacterized than filled by default with the one notion we happen to have. The absence of a rival model is as consistent with that nature being unknown as with the panpsychist's reading.

Russellian monism survives this point, and a Russellian panpsychist may have further arguments. The claim may be defended by appeal to simplicity, explanatory gap considerations, anti-emergentist premises, or a broader theory of categorical properties. But an argument of that form no longer rests on first-person immediacy. Its support is a theory about what best fills an explanatory or metaphysical gap. First-person evidence has unusual authority only within its proper domain. I can be directly aware of my pain. I cannot be directly aware, by that same act, that the intrinsic nature of all matter is pain-like, experience-like, or proto-phenomenal. At most, consciousness supplies a candidate model for intrinsic nature. The suitability of that model must be defended by the theory, since first-person acquaintance supplies the experience itself rather than its fundamental metaphysical role.

\section{The Burden Shift}

Treating panpsychist posits as theoretical rather than first-personal shifts the burden. Micro-experience, proto-experience, or cosmic subjectivity may be possible. The harder question is what these posits explain. Once first-person warrant is no longer available for them, the theoretical phenomenal base must be justified by explanatory work. This burden is not unique to panpsychism; any theory incurs it once its central posits are no longer directly given.

An adequate theory of consciousness should help explain its unity, its contents, its temporal grain, its sensitivity to factors such as sleep and anesthesia, the difference between processes that enter experience and those that do not, and why experience has valence. Panpsychism need not explain all of these features by its fundamental phenomenal base alone. A panpsychist can reasonably say that attention, memory, reportability, temporal integration, and many contents depend on organization. That concession locates the pressure point more precisely.

Organization may explain much of the profile of consciousness as we undergo it. Attention, memory, temporal integration, and reportability may help explain why conscious experience has the structure it has, without explaining why there is anything it is like. Panpsychists can therefore say that the hidden phenomenal base is aimed at phenomenal existence. The theory still has to say what relation that base has to the organized, bounded, valenced consciousness that first-person experience actually gives. If its main role is to block brute emergence from the wholly non-phenomenal, that may be a serious metaphysical role. It is not yet an explanation of the determinate structure of the consciousness we actually undergo. Otherwise the posit marks the need for explanation rather than supplying one. The pressure is sharpest at the combination problem. Micro-experiences were introduced as the right kind of thing to combine into macro-consciousness. If organization does the combining work, the micro-experiences must still make a distinctive contribution in virtue of being experiential. What does the base's being phenomenal explain that a non-phenomenal categorical base, equally available to ground physical structure, could not? If nothing, the same organizational story could in principle be told over non-experiential materials, and the base's phenomenal character does no explanatory work.

This point supplies no complete emergentist theory of consciousness. I am not deriving phenomenality from organization, nor claiming to solve the hard problem; the question remains why any organized processing should feel like anything at all. Still, two explanatory gaps must be distinguished. The first is the physicalist gap: how consciousness arises from a wholly non-conscious base. The second is the gap panpsychism creates for itself: how the unified, determinate consciousness we actually undergo arises from the hidden phenomenal or proto-phenomenal base it posits. These are not the same gap. Panpsychism may have resources against the first, but placing phenomenality at the fundamental level does not close the second, the gap the first-person datum was supposed to have closed in advance. Unless those hidden properties explain enough to justify their theoretical cost, the abductive burden remains in place.

A panpsychist may object that the two gaps are not on a par: the emergence of experience from the wholly non-experiential would be a category error, while the emergence of macro-experience from micro-experience would not. Grant the distinction. It does not restore first-person warrant for the hidden base. Whether macro-from-micro experience escapes a category error is settled, if at all, by an anti-emergentist argument of the kind Section 1 brackets, not by acquaintance; the distinction, if it holds, relocates the defense onto theoretical ground rather than returning it to the immediacy of the datum. And even granted, it does not make the second gap disappear. How micro-experiences combine into the unified, determinate consciousness we undergo remains a substantive task. Panpsychism would have shown that its gap is more tractable in kind, not that the first-person datum has closed it.

Phenomenal-concept strategies press a nearby caution.\footnote{Papineau argues that the felt gap between physical and phenomenal description may arise from how we think about conscious states, and related phenomenal-concept strategies develop the diagnosis in different ways \citep{papineau2002,loar1990,balog2012pcs,carruthersVeillet2007}. Mendelovici's discussion of structural problems across theories of consciousness is another useful comparison \citep{mendelovici2019}.} My claim is narrower than that physicalist strategy. It need not say that phenomenal concepts explain away the hard problem. Even if phenomenal concepts or acquaintance give us a distinctive way of thinking about our own experiences, that fact does not extend first-person warrant to micro-experiences, proto-phenomenal categorical bases, or cosmic subjectivity.

Stoljar's ignorance hypothesis gives a related caution. If part of the pressure toward panpsychism arises from ignorance of the relevant physical or categorical facts, then the proper response is not to deny the datum of consciousness, but to separate the datum from speculative claims about what must underlie it \citep{stoljar2006ignorance}. The introspective warrant constraint is compatible with the ignorance view but does not require it. Whether or not our ignorance of categorical facts is what generates the pressure toward panpsychism, first-person acquaintance is bounded by what it discloses, and that bound holds even if further categorical facts exist. Before adding hidden phenomenal properties to reality, we should ask what explanatory work they perform beyond marking our resistance to brute emergence.

A panpsychist may reply that this burden is shared. Physicalism also owes an account of why organization gives rise to experience, unless it denies or deflates the explanandum, and panpsychism is not alone in facing structural burdens. Independent abductive arguments may also be possible. Panpsychist arguments that make no claim on first-person warrant, of which Saad's harmony argument is the cleanest case, are outside the target. The introspective warrant constraint is a discipline on how first-person evidence is used; it has nothing to say about views that do not use it. First-person evidence cannot by itself pay the cost of the hidden phenomenal base. Where panpsychism is defended by harmony, phenomenal powers, cosmopsychic individuation, or composite subjectivity, its warrant is theoretical and should be evaluated as such.

Dreamless sleep is a useful illustration. It does not prove that consciousness is absent. A panpsychist, a Buddhist philosopher, or a theorist of minimal consciousness may argue that some form of experience occurs without later memory. Such a claim requires independent warrant. If there is no memory, no report, no phenomenological trace, and no distinctive explanatory payoff, then first-person evidence does not license the posit.

Temporal grain illustrates the same pattern. Human consciousness appears to have thresholds and integration windows. Some stimuli may be processed without entering experience at all, and conscious episodes seem to depend on temporally organized uptake rather than on every event in the nervous system surfacing as experience. A panpsychist can say, consistently, that sub-threshold processes still have micro-experiential aspects. But that reply does not yet explain why conscious episodes have the temporal boundaries they have, why some processed information enters experience while other information does not, or why integration at one scale yields a conscious episode while nearby subpersonal processing remains outside awareness. If hidden experientiality is placed beneath the threshold, the question is what it explains about the threshold. Without such an account, the posit multiplies hidden phenomenal facts beneath temporal structure rather than explaining that structure.

These examples illustrate the warrant problem; they are not empirical proofs against panpsychism. Panpsychism repeatedly posits experience or experience-like properties where experience itself gives no clear warrant. That may be permissible as theory. It is not directly licensed by the first-person datum. Worries about introspection strengthen the restriction. Schwitzgebel argues that introspection may be unreliable even in cases where we might expect it to be secure \citep{schwitzgebel2008,schwitzgebel2011}. If introspection is fallible even within its proper domain, then its extension to hidden micro-, proto-, or cosmic phenomenal structure requires special caution.

\section{Conclusion: The Limits of Special Warrant}

First-person experience gives panpsychism a strong beginning. It tells us that consciousness is real, and something about what conscious life is like; that is real evidence, and the paper does not minimize it. The same evidence does not tell us that consciousness is fundamental, microphysical, cosmic, or proto-phenomenal. Panpsychism's further phenomenal posits therefore need a different warrant. They may still be defended by the explanatory work they do, but knowing consciousness from within does not by itself settle its metaphysical status.

Panpsychism's central promise changes accordingly. The view appeared to dissolve the explanatory gap in advance by placing phenomenality at the fundamental level, so that consciousness as we undergo it did not have to arise from anything wholly non-phenomenal. Once first-person inheritance is removed and the hidden phenomenal level is admitted as a theoretical posit, that dissolution is no longer secured in advance. The gap reappears between the hidden base and the consciousness we actually undergo. Goff's recent treatment of subjects as strongly emergent while phenomenal properties are weakly emergent makes the reappearance visible within his own framework \citep{goff2024}. Panpsychism is entitled to try to close the gap. It is not entitled to claim that the first-person datum has already closed it.

Panpsychism may appeal to independent metaphysical arguments, anti-emergentist principles, simplicity considerations, psychophysical harmony, or Russellian claims about intrinsic nature. Those arguments may succeed or fail on their merits, but they may not borrow the authority of first-person experience and spend it on entities or structures that first-person experience never reveals. Once the first-person route gives way to abduction, austere realistic monism, micropsychism, cosmopsychism, panprotopsychism, and phenomenal-powers grounding must compete with their rivals on the same abductive terms as any other metaphysical hypothesis.

Nothing here vindicates illusionism or deflationary physicalism, since those views face their own burden: denying or deflating the first-person datum carries a substantial cost. Any view that takes first-person consciousness seriously must also respect the limits of what first-person consciousness discloses.

Where panpsychism treats the epistemic privilege of first-person experience as support for micro-experience, proto-phenomenal bases, cosmic subjectivity, or other hidden phenomenal posits, it overdraws on experience. Panpsychism in its first-person-motivated forms is the central case considered here, but the same constraint applies more generally. Panpsychism can still be defended, but not as the view whose metaphysical base is already given in experience. That base is hidden, theoretical, and abductive. Panpsychism's advantage over rival theories must therefore be earned, not inherited.
